\documentclass{article}

\usepackage{amsmath,amssymb}
\usepackage{xurl}
\usepackage[hidelinks]{hyperref}
\setlength{\emergencystretch}{2em}

% Shared commands for notes in docs/paper-gaps/.

\DeclareUnicodeCharacter{2080}{\ifmmode{}_{0}\else\textsubscript{0}\fi}
\DeclareUnicodeCharacter{2081}{\ifmmode{}_{1}\else\textsubscript{1}\fi}
\DeclareUnicodeCharacter{2082}{\ifmmode{}_{2}\else\textsubscript{2}\fi}
\DeclareUnicodeCharacter{2099}{\ifmmode{}_{n}\else\textsubscript{n}\fi}

\newcommand{\Lean}{\textsc{Lean}}
\ExplSyntaxOn
\NewDocumentCommand{\leanid}{m}{
  \ifmmode
    \textsf{\detokenize{#1}}
  \else
    \group_begin:
    \urlstyle{sf}
    \tl_set:Nn \l_tmpa_tl {#1}
    \tl_replace_all:Nnn \l_tmpa_tl {\_} {_}
    \exp_args:NV \path \l_tmpa_tl
    \group_end:
  \fi
}
\ExplSyntaxOff
\providecommand{\tr}{\operatorname{tr}}
\newcommand{\tnleanIssueBase}{https://github.com/LionSR/TNLean/issues/}
\newcommand{\tnleanPrBase}{https://github.com/LionSR/TNLean/pull/}
\newcommand{\tnleanDocBase}{https://sirui-lu.com/TNLean/docs/}
\newcommand{\ghissue}[1]{\href{\tnleanIssueBase#1}{GitHub issue \##1}}
\newcommand{\ghpr}[1]{\href{\tnleanPrBase#1}{PR \##1}}
\newcommand{\doclink}[1]{\href{\tnleanDocBase#1.html}{\path{#1}}}

% Dirac notation and the identity operator, as in blueprint/src/macros/common.tex.
% \providecommand keeps note-local definitions authoritative.
\usepackage{dsfont}
\providecommand{\ket}[1]{|#1\rangle}
\providecommand{\bra}[1]{\langle #1|}
\providecommand{\braket}[2]{\langle #1 | #2 \rangle}
\providecommand{\ketbra}[2]{|#1\rangle\!\langle #2|}
\providecommand{\Id}{\mathds{1}}
\providecommand{\one}{\mathds{1}}
\providecommand{\C}{\mathbb{C}}
\providecommand{\MN}[1]{M_{#1}(\C)}
\providecommand{\MD}{M_D(\C)}

% External referencing for paper sources.  The build helper writes sanitized
% auxiliary files under build/paper-aux/ and excludes duplicated source labels.
\usepackage{xr}

% Import a generated paper auxiliary file with a caller-supplied label prefix.
% The candidate paths support builds from docs/paper-gaps/, the repository root,
% and build/paper-aux/ itself.
\newcommand{\importpaperaux}[2]{%
  \IfFileExists{../../build/paper-aux/#2.aux}{%
    \externaldocument[#1]{../../build/paper-aux/#2}%
  }{%
    \IfFileExists{build/paper-aux/#2.aux}{%
      \externaldocument[#1]{build/paper-aux/#2}%
    }{%
      \IfFileExists{#2.aux}{%
        \externaldocument[#1]{#2}%
      }{}%
    }%
  }%
}

% General external reference macros
\newcommand{\paperref}[2]{\ref{#1-#2}}
\newcommand{\papereqref}[2]{(\ref{#1-#2})}

% CPSV16 source (1606.00608 / MPDO-22-12-17-2.tex) external document and shortcuts
\importpaperaux{cpsv16-}{1606.00608/MPDO-22-12-17-2-tnlean}
\newcommand{\cpsvref}[1]{\paperref{cpsv16}{#1}}
\newcommand{\cpsveqref}[1]{\papereqref{cpsv16}{#1}}

% Verdict marker: \gapnote{<kind>}{<status>}. Kind is the policy
% classification (clarification, local-correction, scope-restriction,
% unfaithful, false-source, open-gap); status is open, wip (elimination
% underway), resolved, or historical. Machine-read by the site index;
% prints a verdict line.
\newcommand{\gapnote}[2]{\par\noindent\textbf{Verdict:} #1 (#2).\par}


\title{Dependency Audit for the Non-periodic MPS Fundamental Theorem Blueprint}
\author{}
\date{2026-05-08}

\begin{document}
\maketitle
\gapnote{clarification}{open}

\section{Scope}

This note answers the organizational questions in
\href{https://github.com/LionSR/TNLean/issues/1530}{issue~1530} by comparing
the blueprint dependencies with the non-periodic Fundamental Theorem route in
Refs.~\cite{Cirac2016MPDO_arXiv,Cirac2021Matrix}. It does not address the
direct periodic theorem of Ref.~\cite{DeLasCuevas2017Irreducible} or the
fixed-length algebraic theorem of Ref.~\cite{Molnar2018NormalPEPS}. It also
identifies which hypotheses in the current conditional after-blocking
statements are outputs expected from the source route and which remain
unproved inputs in the blueprint.

The relevant source passages are the basis-of-normal-tensors definition,
Proposition~\texttt{prop:char-BNT},
Equation~\texttt{eq:II\_ABasicTensors}, the block-injective input
Proposition~\texttt{propblockinj}, the proportional Fundamental Theorem
\texttt{thm1}, and the equal-MPV Corollary~\texttt{II\_cor2} in
Ref.~\cite[lines~271--356]{Cirac2016MPDO_arXiv}. In the appendix of the same
source, Lemma~\texttt{Lem:app\_simple} appears at lines~1155--1163,
Theorem~\texttt{thm1} is restated at lines~1167--1170 and proved at
line~1182, and Corollary~\texttt{II\_cor2} appears at lines~1172--1179.
Equation~\texttt{eq:II\_auxcor} of
Ref.~\cite[lines~1175--1192]{Cirac2016MPDO_arXiv} states the global gauge
formula and derives it in lines~1189--1192. Section~IV of
Ref.~\cite{Cirac2021Matrix} gives the corresponding review treatment. These
source theorems provide the standard of comparison. Formal helper structures
are compatible with that standard only when the blueprint identifies them as
helpers rather than source hypotheses or conclusions.

\paragraph{Notation.}
Let $\mathbb{C}$ be the field of complex numbers, let
$\mathbb{C}^{\times}$ be its multiplicative group, and let
$M_D(\mathbb{C})$ denote the algebra of complex $D\times D$ matrices. Let $d$
be the physical dimension of a matrix product state tensor
$A=(A^i)_{i=1}^d$. For a positive integer $p$, let $A^{[p]}$ denote the
length-$p$ blocking of $A$; its physical dimension is $d^p$. For a word
$w=i_1\cdots i_L$, define its length by $|w|=L$ and its associated product by
$A^w=A^{i_1}\cdots A^{i_L}$.

A basis of normal tensors, abbreviated BNT, is a family
$(A_j)_{j=1}^g$ of representative normal tensors. For such a family, define
the block-diagonal tensor $\tilde A$ by
\begin{align}
    \tilde A^i
        &= \bigoplus_{j=1}^g A_j^i,
    \label{eq:tnlean_ft_dependency_audit_bnt_direct_sum}\\
    \tilde A^w
        &= \tilde A^{i_1}\cdots\tilde A^{i_L}.
    \label{eq:tnlean_ft_dependency_audit_bnt_word_product}
\end{align}
The abbreviation MPV denotes the matrix product vector generated by a tensor.
Let $\mathbf{0}_{d^p,z}$ denote the zero tensor with physical dimension $d^p$
and bond dimension $z$, and let $A_{\mathrm{nz}}$ denote the nonzero part of a
tensor after removal of a zero direct summand. The relation $\sim$ denotes the
canonical-form direct-sum decomposition up to the usual gauge change.

For a BNT comparison, let $\pi$ be a permutation of the representative normal
tensors, let $X_j$ be an invertible gauge matrix, and let
$\zeta_j\in\mathbb{C}^{\times}$ be the nonzero phase factor in the resulting
relation between matched representatives.

\section{Block permutation and separation}

The blueprint chapter ``Block Permutation and Separation'' contains several
distinct ingredients. Its product-algebra argument considers the algebra
\begin{align}
    \mathcal A
        &= \prod_{j=1}^g M_{D_j}(\mathbb{C}).
    \label{eq:tnlean_ft_dependency_audit_product_algebra}
\end{align}
An automorphism of $\mathcal A$ permutes its simple factors and acts by an inner
automorphism on each factor. This algebraic mechanism supports the
same-structure and block-injective comparison statements. The chapter's
invariant-projection results instead belong to canonical-form reduction. Its
block-separation lemmas show that, under strict canonical-form hypotheses, an
identity between the MPV families of two block-diagonal tensors separates into
identities between the MPV families of their individual blocks.

The current conditional non-periodic Fundamental Theorem does not depend on
this chapter as a single unit. Its visible downstream dependencies are the
canonical-form condition imposed on BNT data and the same-structure comparison
lemma. That lemma treats the special case in which the tensors already have
the same number of blocks, the same block dimensions, and the same weights. It
does not establish the general BNT matching step of
Ref.~\cite{Cirac2016MPDO_arXiv}.

The general source proof instead uses BNT representatives, eventual linear
independence, a block-injective fixed-length span, and scalar power sums, as
shown in Ref.~\cite{Cirac2016MPDO_arXiv}. After blocking, the required
homogeneous block-injective direct-sum span is
\begin{align}
    \operatorname{span}
        \left\{\tilde A^w:|w|=L\right\}
        &= \bigoplus_{j=1}^g M_{D_j}(\mathbb{C}).
    \label{eq:tnlean_ft_dependency_audit_block_injective_span}
\end{align}
An equivalent trace-dual separation statement may replace
Eq.~\ref{eq:tnlean_ft_dependency_audit_block_injective_span}. This input
remains missing from the general blueprint route. It is recorded separately
in \path{docs/paper-gaps/rmp_nonperiodic_bnt_comparison_inputs.tex} and
\path{docs/paper-gaps/pgvwc07_direct_sum_input.tex}.

The block-permutation chapter should therefore remain before the non-periodic
Fundamental Theorem. The blueprint should not, however, present every result in
that chapter as a dependency of the remaining BNT comparison argument.

\section{Multiplicative domain versus operator convexity}

The Kadison--Schwarz and multiplicative-domain part of the Schwarz chapter is a
genuine dependency of the aperiodic route. For normalized completely positive
maps, the argument places peripheral eigenvectors in the multiplicative domain
and then derives the commutation and phase-gauge rigidity used in the spectral
and canonical-form chapters; see Ref.~\cite{Wolf2012Quantum} for the underlying
quantum-channel framework. This material must remain before Perron--Frobenius
theory and the non-periodic Fundamental Theorem of
Refs.~\cite{Cirac2016MPDO_arXiv,Cirac2021Matrix}.

The material on operator convexity, Jensen inequalities, and Lieb concavity has
a different role. An audit of the current blueprint labels finds no use of that
section in the non-periodic Fundamental Theorem chapter or its BNT comparison
chain. Its natural applications lie in later analytic arguments concerning
entropy, matrix product density operators, and parent Hamiltonians; related
entropy applications appear in
Refs.~\cite{Beigi2013Sandwiched,MullerLennert2013Renyi}. In particular, this
material is relevant to the open operator-convexity and POVM-resolution work
around the parent-Hamiltonian gap.

Moving the entire Schwarz chapter would therefore break a real dependency. A
safe later reorganization would retain Kadison--Schwarz and the multiplicative
domain upstream while moving only the operator-convexity and Jensen material
to an analytic chapter after the Fundamental Theorem.

\section{Which assumptions are extra}

Some hypotheses in the conditional after-blocking statements are expected
outputs of the source's canonical-form reduction rather than additional
mathematical assumptions. They include a common positive blocking length,
trace-preserving, primitive, and irreducible nonzero sectors, and the
zero-summand decompositions
\begin{align}
    A^{[p]}
        &\sim \mathbf{0}_{d^p,z_A}\oplus A_{\mathrm{nz}},
    \label{eq:tnlean_ft_dependency_audit_zero_tail_a}\\
    B^{[p]}
        &\sim \mathbf{0}_{d^p,z_B}\oplus B_{\mathrm{nz}},
    \label{eq:tnlean_ft_dependency_audit_zero_tail_b}\\
    z_A
        &= z_B.
    \label{eq:tnlean_ft_dependency_audit_zero_tail_dimension}
\end{align}
Once the upstream canonical-form and period-removal arguments have been
proved, the source route supplies the data in
Eqs.~\ref{eq:tnlean_ft_dependency_audit_zero_tail_a}--%
\ref{eq:tnlean_ft_dependency_audit_zero_tail_dimension}
\cite{Cirac2016MPDO_arXiv,Cirac2021Matrix}.

Other hypotheses remain genuine proof obligations in the current conditional
statements. These are one-site injectivity at the selected blocking length and
the BNT block-matching input required to compare representatives. The source
derives this matching from the proportional-MPV hypothesis and the
canonical-form BNT structure \cite{Cirac2016MPDO_arXiv}. The required
conclusion consists of a permutation $\pi$, equality of the matched bond
dimensions, invertible gauges $X_j$, and nonzero phase factors
$\zeta_j\in\mathbb{C}^{\times}$ such that
\begin{align}
    B_{\pi(j)}^i
        &= \zeta_j X_j A_j^i X_j^{-1}.
    \label{eq:tnlean_ft_dependency_audit_bnt_block_matching}
\end{align}

The former formal route through coefficient arrays with nonzero limiting
coefficients has been deleted because it does not match the source theorem's
statement. The remaining formal statements assume
Eq.~\ref{eq:tnlean_ft_dependency_audit_bnt_block_matching}; they do not derive
it from equality of the MPV families. Those statements only transport the
assumed matching through the common blocking and recover the sector weights.

The blueprint and the associated paper-gap notes must therefore continue to
mark this passage as conditional. This audit does not claim to resolve
issues~1499, 1500, or 1501.

\section{Staged blueprint organization}

The chapter order can be improved without changing any theorem statement. The
aperiodic dependencies should appear first and end with the non-periodic
Fundamental Theorem of
Refs.~\cite{Cirac2016MPDO_arXiv,Cirac2021Matrix}. The direct periodic theorem
of Ref.~\cite{DeLasCuevas2017Irreducible}, the material on symmetries and
string order, and the fixed-length algebraic theorem of
Ref.~\cite{Molnar2018NormalPEPS} should form the first block after the
Fundamental Theorem. Entropy and matrix product density operator material
should precede parent Hamiltonians because the parent-Hamiltonian chapter
depends conceptually on state and density-operator methods, not conversely.

A larger split of the Schwarz chapter should occur in a separate change. Once
the operator-convexity section has been isolated, that split is purely
organizational, but it affects many labels and should not be combined with
claims that any proof gap has been closed.

\section{Conclusion}

The dependency audit separates established upstream requirements from
unproved comparison inputs. Kadison--Schwarz, the multiplicative domain,
canonical-form reduction, and the relevant block-permutation results remain
upstream dependencies. Operator-convexity and Jensen material may move
downstream. The general non-periodic route remains conditional on the
block-injective span in
Eq.~\ref{eq:tnlean_ft_dependency_audit_block_injective_span} and the BNT
matching conclusion in
Eq.~\ref{eq:tnlean_ft_dependency_audit_bnt_block_matching}. No reorganization
described here resolves the outstanding proof obligations.

\bibliographystyle{alpha}
\bibliography{references}

\end{document}
